\section{References}
\label{sec:regular-bib}

The name ``regular'' for the functions studied in this part comes from~\cite[p. 217]{engelfrietHoogeboom2001}. Although we begin with a definition in terms of prime functions, this is a later development historically, and the   original definition was in terms of two-way transducers.  The decomposition into primes is implicit in~\cite{bojanczykDaviaudKrishna2018} and explicit in~\cite[Theorem 18]{bojanczykStefanski2020}, although the latter paper discuss a more general framework of orbit-finite sets. 

\paragraph*{Two-way transducers.}
Two-way automata as acceptors (and not transducers) appear in \cite[Definition 2]{shepherdson1959reduction} and \cite[Definition 13]{RabinScott59}. Both of these papers prove that such acceptors define exactly the regular languages, which covers the continuity result from \cref{thm:continuity-2dfas}. Shepherdson also sketches a transducer variant of two-way automata, providing the first known appearance of the two-way transducer model that we use here~\cite[Note 4]{shepherdson1959reduction}. The transducer model is also somewhat recognisable in papers from the 1960s and 1970s such as~\cite{hopcroftUllman1967} or~\cite{AhoHopcroftUllman1969}. In the latter case, our two-way transducer model is the same as balloon automata with no auxilliary storage. 


Closure of two-way transducers under pre-composition with Mealy machines was first proved in~\cite[Theorem 5]{hopcroftUllman1967} using an ingenious construction. In our proof the same result, see~\cref{lem:2dfa-precomposition-with-mealy}, we avoid ingenuity by appealing to the decompositions into prime functions. 
Closure of two-way transducers under composition, which is \cref{thm:composition-of-two-way-transducers} in this book,  was first proved in \cite[Theorem 2]{ChytilJakl1977}. The authors of \cite{ChytilJakl1977} credit this result to~\cite{AhoUllman1970}. I believe that they are referring to \cite[Theorem 2]{AhoUllman1970}, but that theorem talks about inverse images of languages under two-way automata, and not about compositions of two transducers, although the proof does seem to contain the right ingredients for composition of transducers.


\paragraph*{Decidability of equivalence.}
Decidability of equivalence for two-way transducers was first proved in~\cite[Theorem 1]{gurari1982}. Our statement from \cref{thm:decidable-equivalence-regular} deviates from that result in that it uses a different model (compositions of primes) which is only later proved to be equivalent,  and a proof that is based on weighted automata,  see the comments in~\cref{sec:rational-bib}.


\paragraph*{Streaming string transducers.} This model was introduced in~\cite[Section 3]{alurCerny2010}, where it was shown to be equivalent to two-way transducers and \mso transductions in Theorems 1,2,3. Many of the main ideas behind this model can be traced back to~\cite[Theorem 17]{bloemEngelfriet2000}. If we restrict the result from \cite{bloemEngelfriet2000} from trees to strings, and translate it to our terminology, then it says that the \mso transductions are exactly those that can be computed by an \sst with regular look-ahead (regular look-ahead can be formalised as a pre-composition with a right-to-left Mealy machine). Now it is known that the regular look-ahead can be eliminated, since already without look-ahead \sst's are equivalent to the regular functions.  Another early transducer take on the regular functions (again, using \mso transductions and the more general tree-to-tree case) can be found in~\cite[Theorem 7.1]{engelfrietManeth1999}; the paper uses macro tree transducers, which are nicely explained in~\cite{nguyen2024} using a terminology similar to the one of this book. Despite the earlier work described above, one should note that the attractive operational presentation of \sst's from~\cite{alurCerny2010} has proved instrumental in attracting attention to the topic.


\paragraph*{Logic.}  The original references for \cref{thm:mso-logic-languages} about the equivalence of \mso and regular languages are  \cite[Theorems 1 and 2]{buchi1960},  \cite[Theorem on p. 35]{elgot1961} and \cite[\rus{Теорема 10 и 3}]{trakhtenbrot1962}. Interestingly enough, the logical corollary that was the initial goal of this work, which is that weak \mso is decidable on the natural numbers with order, is credited to unpublished -- and possibly automata-free -- work of Ehrenfeucht in~\cite[Section 5.8]{elgot1961}. The connections between \mso and automata are a big topic, a good introduction is the survey paper \cite{thomas1997}. 

A version of \cref{thm:logic-rational-functions} about \mso relabelings can be found in~\cite[Theorem 10]{bloemEngelfriet2000}, with some minor differences: it is letter-to-letter and uses trees instead of strings. The logical characterisation of the regular functions from \cref{thm:logic-regular-functions} is due to \cite{engelfrietHoogeboom2001}, where it was shown that \mso transductions have the same expressive power as two-way transducers. In my opinion, this is a foundational paper for the regular functions, since it shows that they can be descibed in at least two different ways, including  logic (also, this paper establishes the name ``regular'' for this class.

The automata characterisation of the first-order fragment that is discussed in \cref{sec:logic-aperiodic} is a combination of two results: a breakthrough result of Sch\"utzenberger which showed that star-free regular expressions are the same as aperidodic monoids \cite[Main Property]{Schutzenberger1965}, and an important later addition of McNaughton and Papert which showed that star-free regular expressions are the same as first-order logic \cite[Theorem 10.5]{McNaughtonPapert1971}. Our proof leverages the Krohn-Rhodes Theorem, this option was already noticed in~\cite[Footnote 1]{meyer1969}.

\paragraph*{Combinators.} \cref{thm:regular-terms} was originally proved in \cite[Theorem 6.1]{bojanczykDaviaudKrishna2018}. There is an alternative approach to regular expressions for transducers, which does not use types~\cite[Theorem 15]{alurFreilichRaghothaman2014}.